\subsection{Sources of Dynamic Imbalance}
\label{section:dynamic-imbalance}

Due to the heterogeneity of multimodal model architectures and the diversity of training data,
the training process for LMMs differs significantly from that for unimodal models in
\emph{pipeline stage imbalance} and \emph{training data dynamicity}.

\heading{Pipeline Stage Imbalance}
Optimal pipeline efficiency requires balanced stage partitioning to minimize bubbles, yet inherent computational disparities between modality modules introduce fundamental challenges.
Consider a 37B VLM comprising a 5B ViT encoder (64 layers) and 32B language model (64 layers) on H800 GPUs,
processing a batch with 8 images and 8192 text tokens.
Each ViT layer processes the images in 6.75ms (forward+backward), while each LM layer requires 10.5ms for the tokens.
The optimal partitioning scheme produced by exhaustive search across 16 pipeline stages yields stage latencies between 63ms and 73.5ms (\ie 16.7\% variation).
With 64 microbatches and Megatron-LM's 1F1B pipeline schedule, this imbalance introduces 22.8\% additional pipeline bubbles, demonstrating the inherent difficulty in achieving perfectly balanced partitioning.

\heading{Training Data Dynamicity}
This challenge is further compounded by the dynamic heterogeneity in multimodal training data,
where the variability across batches from different modalities induces fluctuating computational demands.
Multimodal data originates from diverse sources including:
image-description pairs~\cite{laion5b-nips22,coyo-700m},
video with captions~\cite{share-gpt4-video-arxiv24,internvid-arxiv23,mmtrail-arxiv24,msr-vtt-cvpr16}, and
interleaved image-text content~\cite{obelics-arxiv23,mmc4-arxiv23}.
Data samples typically comprise images/videos with descriptive texts, exhibiting significant \emph{cross-dataset variation} in modality data ratios (\figref{figure:data-distribution}a--b).
For instance,
the LAION-2B~\cite{laion5b-nips22} dataset consists of images paired with short captions, with a small text-image ratio (16.4 tokens/image), while the OBELICS~\cite{obelics-arxiv23} dataset contains full-length multimodal documents
with highly variable ratios (0.4 to 3115 tokens/image).

This variation further induces \emph{cross-batch imbalance} during data packing.
Typically, data samples are packed into larger batches to enhance computation efficiency~\cite{dynapipe-eurosys24,wlb-llm-osdi25}.
In vision-language models (VLMs),
both images and texts are tokenized (\eg one image into 169 patch tokens)
and then greedily packed up to the model's context length (\eg 8192 tokens) to form a microbatch.
Similarly, for text-to-video (T2V) diffusion models,
video clips with similar durations and aspect ratios are often grouped for batch processing~\cite{moviegen-arxiv24,step-t2v-arxiv25}.
However, data packing fails to ensure balanced workloads across multiple modalities
due to their divergent distributions (\figref{figure:data-distribution}c--d).
For instance, consider two microbatches with same 8192-token capacity:
the first contains 10 images (\ie 1690 patch tokens) accompanied by 6502 text tokens,
whereas the second consists of only a single image and 8023 text tokens.
These two microbatches impose disparate computational demands:
the FLOPs required for the image pipeline stages in the first microbatch are $10\times$ higher than those in the second.
This discrepancy leads to imbalanced pipeline stages across microbatches (\figref{figure:imbalanced-pipeline}).
Furthermore, for a T2V model comprising a 7B LM and a 5B DiT,
the most computationally intensive batch (data batch index 100 in \figref{figure:data-distribution}d) imposes a $4.15\times$ greater computational load than the smallest batch (data batch index 1),
even after data packing.

% This variation further induces \emph{cross-batch imbalance} during data packing.
% Data samples can be packed into larger batches for better computation efficiency~\cite{dynapipe-eurosys24,wlb-llm-osdi25}.
% For vision-language models (VLMs), both images and texts are first encoded into tokens (\eg one image into 169 patch tokens).
% After that, data samples are greedily packed up to the context length of the model (\eg 8192 tokens) to form one microbatch.
% For text-to-video (T2V) diffusion models, video clips with similar durations and aspect ratios can be grouped together and processed in one batch~\cite{moviegen-arxiv24,step-t2v-arxiv25}.
% However, data packing fails to achieve balanced workloads when dealing with multiple modalities due to their divergent distributions (\figref{figure:data-distribution}c--d).
% For example, suppose there are two microbatches both with 8192 tokens, the first microbatch contains 10 images (\ie 1690 patch tokens) and 6502 text tokens,
% while the second microbatch may only contain text tokens.
% These two microbatches will exhibit very distinct computation requirements and lead to imbalanced pipeline stages,
% For a T2V model with 7B LM and 5B DiT, the largest data sample (data batch index 100 in \figref{figure:data-distribution}d) still incurs $4.15\times$ greater computational load than the smallest (data batch index 1) after data packing.

\subsection{Negative Impact on Pipeline Parallelism}
\label{section:impact-on-pipeline-parallelism}

The combination of pipeline stage imbalance and training data dynamicity leads to the \emph{dynamic imbalance} problem.
As depicted in \figref{figure:imbalanced-pipeline}, this issue introduces significant pipeline bubbles in Megatron-LM's 1F1B pipeline schedule, and severely degrades training throughput.

To quantify this impact, we compare two model setups with the same number of parameters:
a unimodal LM (7B) and
a vision-language model (ViT 2B + LM 5B).
Under identical 1F1B pipeline configurations
(using balanced parameter partitioning and fixed computational budgets),
the VLM incurs a 12.5\% overhead on static data due to stage imbalance.
With real-world dynamic data, this overhead escalates to 40.3\%,
as shown in \tableref{table:imbalanced-pipeline}.

\begin{table}[t]
    \caption{
        Training performance of 7B models on 8 GPUs (TP=2, PP=4).
        ``PFLOPs'' denotes the floating-point operations per iteration in petaFLOPs,
        which is controlled to a fixed value for fair comparison.
        ``MFU'' refers to model FLOPs utilization.
    }
    \label{table:imbalanced-pipeline}
    \small
    \begin{tabular}{>{\raggedright}p{4cm}ccc}
    \toprule
    \textbf{Model Setup} & \textbf{Time (s)} & \textbf{PFLOPs} & \textbf{MFU} \\
    \midrule
    LM 7B & 4.068 & 12.8 & 0.400 \\
    ViT 2B + LM 5B (static data) & 4.567 & 12.7 & 0.351 \\
    ViT 2B + LM 5B (dynamic data) & 6.789 & 12.8 & 0.239 \\
    \bottomrule
    \end{tabular}
\end{table}

\heading{Impact on Pipeline Design}
Such performance degradation renders the fixed pipeline schedules used in unimodal LLMs impractical.
Since each iteration processes a distinct data batch, the optimal pipeline schedule changes dynamically.

Precomputation of pipeline schedules is also infeasible. The number of possible input combinations grows exponentially in two dimensions: the number of modalities and the number of microbatches. In our experimental setup, each microbatch can contain up to 48 images (\S\ref{section:datasets}), leading to $49^{64} \approx 1.5\mathrm{e}108$ distinct input configurations for 64 microbatches. This makes it impossible to exhaustively precompute optimal schedules for all scenarios. Although precomputing over a subset of inputs could reduce overhead, selecting a representative subset from this enormous space is highly challenging, and fails to generalize to unseen inputs.
% Moreover, any changes to the dataset (\eg new data or reshuffling) or parallelism configurations (\eg TP/PP sizes) also invalidate precomputed schedules, necessitating costly regeneration.

Several prior works have addressed the issue of dynamicity in text sequence lengths,
by improving data packing strategies to reduce discrepancies between microbatches~\cite{wlb-llm-osdi25},
or adaptively reordering microbatches to minimize bubbles caused by irregular inputs~\cite{dynapipe-eurosys24,flexpipe-atc25}.
However, these methods are inadequate for solving dynamic imbalance in multimodal training, as they overlook the heterogeneity among modality modules.
Unlike in unimodal LLMs, where sequence length variations affect all layers uniformly, multimodal microbatches affect modality-specific modules differently. For example, increasing the number of images substantially raises computational demands on image encoders, while barely affecting the language model backbone. This inter-modality imbalance makes data-centric approaches designed for unimodal models ineffective for LMM training.
